\def\Ponderation{40}
\def\SigleNumero{STT-1000}
\def\NomCours{Probabilités et statistique}
\def\DateEvaluation{Hiver 2026}
\def\HeureEvaluation{}
\def\Session{}
\def\NomEvaluation{Examen 2 - Reprise}

% Ne rien mettre entre les accolades si l'enseignant (ou l'enseignante) est aussi le coordonnateur (ou la coordonnatrice)
\def\Coordonnatrice{}
\def\Coordonnateur{}

% Ne rien mettre entre les accolades s'il y a plus d'une section
\def\Enseignant{}
\def\Enseignante{}

% Ne rien mettre entre les accolades s'il s'agit d'un cours à section unique
% Mettre le nom de chacun des enseignants et enseignantes entre les accolades
% Une case à cocher sera générée pour chacune des sections
\def\SectionA{Anne-Sophie Charest}
\def\SectionB{Julien Miron}
\def\SectionC{}
\def\SectionD{}


\def\Directives{
\begin{enumerate}
\item Matériel autorisé:
\begin{itemize}
\item Calculatrice scientifique {\bf autorisée par la faculté des sciences et de génie};
\item Le feuillet de formules fourni avec l'examen. Celui-ci inclut la table de la loi normale. Veuillez ne rien écrire sur le feuillet de formules.
\end{itemize}
\item Assurez-vous que cette évaluation comporte \numquestions ~exercices répartis sur \numpages{}~pages, pour un total de \numpoints{}~points.\
\item \textbf{Veuillez ne pas dégrafer votre examen.}
\item Vous devez\textbf{ justifier toutes vos réponses} par une démarche claire, sauf pour les questions à choix de réponse.\
\item Pour les questions à choix de réponse et les vrais ou faux, vous pouvez mettre un X, un crochet ($\checkmark$) ou remplir la case.\
\item Vous devez laisser tout ce qui ne servira pas durant l'examen à l'avant de la salle AVANT de rejoindre votre place.\
\item Vous n'avez droit à AUCUN appareil électronique en votre possession durant l'examen.
\item Une fois assis à votre place, vous devez demandez l'autorisation pour vous lever, même si l'examen n'est pas encore commencé, sauf lorsque vous avez terminé l'examen.
\item Dans les 5 dernières minutes de l'examen, il sera interdit de se lever avant que TOUTES les copies n'aient été ramassées.
\end{enumerate}
\begin{itemize}
\item[$\rightarrow$] Le non-respect des consignes 7-8-9 entraînera une pénalité de 10\%.
\end{itemize}
Attention, la formule correcte pour les degrés de liberté $\nu$ est 
$$\nu = \frac{(s_1^2/n_1 + s^2_2/n_2)^2}{ \frac{(s_1^2/n_1)^2}{n_1 +1} + \frac{(s^2_2/n_2)^2}{n_2 + 1} } -2 $$
}

\providecommand{\Gradescope}{}


\documentclass{Evaluation}
\usepackage{multicol,graphicx}

\noprintanswers
\begin{document}

\begin{questions}


\pageblanche

%%%%%%%%%%%%%%%% QUESTIONS RÉCAPITULATIFS - MODULES 1 À 5 %%%%%%%%%%%%%%%%

\question Dites si les énoncés suivants sont vrais ou faux. Aucune justification n'est requise.

\begin{parts}

    \part[1] Soit $x_1,\ldots,x_n$. Si la médiane vaut $25$, environ $50\%$ des observations sont supérieures à $25$ et environ $50\%$ sont inférieures à $25$.

    \begin{oneparcheckboxes}
        \correctchoice Vrai
        \choice Faux
    \end{oneparcheckboxes}

    \begin{solution}
    Vrai. C'est la définition de la médiane : elle partage l'échantillon en deux moitiés égales.
    \end{solution}

    \part[1] Soit $x_1,\ldots,x_n$. Si on ajoute une constante $c>0$ à chaque observation, l'écart-type ne change pas.

    \begin{oneparcheckboxes}
        \correctchoice Vrai
        \choice Faux
    \end{oneparcheckboxes}

    \begin{solution}
    Vrai. L'écart-type mesure la dispersion autour de la moyenne. Une translation de toutes les observations ne modifie pas les distances entre les valeurs et leur moyenne.
    \end{solution}

    \part[1] Soit $x_1,\ldots,x_n$. L'étendue est une mesure de dispersion résistante à la présence de valeurs aberrantes.

    \begin{oneparcheckboxes}
        \choice Vrai
        \correctchoice Faux
    \end{oneparcheckboxes}

    \begin{solution}
    Faux. L'étendue est définie comme la différence entre le maximum et le minimum, et est donc très sensible aux valeurs extrêmes.
    \end{solution}

    \part[1] Soit $x_1,\ldots,x_n$. Si la variance échantillonnale $S^2$ vaut $0$, alors toutes les observations sont identiques.

    \begin{oneparcheckboxes}
        \correctchoice Vrai
        \choice Faux
    \end{oneparcheckboxes}

    \begin{solution}
    Vrai. $S^2=0$ implique que chaque terme $(X_i-\overline{X})^2=0$, donc $X_i=\overline{X}$ pour tout $i$.
    \end{solution}

\end{parts}

\vspace{0.5cm}

\question[4]
Soit deux événements $A$ et $B$ formant une partition de l'espace échantillonnal, avec $P(A)=0{,}6$. On sait aussi que $P(D\mid A)=0{,}05$ et $P(D\mid B)=0{,}10$.

Quelle expression permet de calculer $P(A\mid D)$?\\

\begin{checkboxes}
\choice $\dfrac{0{,}4\times 0{,}10}{0{,}05+0{,}10}$\\[4pt]
\choice $\dfrac{0{,}05}{0{,}6\times 0{,}05+0{,}4\times 0{,}10}$\\[4pt]
\correctchoice $\dfrac{0{,}6\times 0{,}05}{0{,}6\times 0{,}05+0{,}4\times 0{,}10}$\\[4pt]
\choice $\dfrac{0{,}4\times 0{,}10}{0{,}6\times 0{,}05+0{,}4\times 0{,}10}$\\[4pt]
\choice $\dfrac{0{,}6\times 0{,}05}{0{,}6\times 0{,}05+0{,}6\times 0{,}10}$
\end{checkboxes}

\begin{solution}
La bonne réponse est \textbf{c)}.

D'après le théorème de Bayes :
$$P(A\mid D)=\frac{P(D\mid A)\,P(A)}{P(D\mid A)\,P(A)+P(D\mid B)\,P(B)}=\frac{0{,}6\times 0{,}05}{0{,}6\times 0{,}05+0{,}4\times 0{,}10}.$$

Problèmes avec les autres réponses :
\begin{itemize}
\item \textbf{a)} utilise $P(B)=0{,}4$ au numérateur avec $P(D\mid B)=0{,}10$, au lieu de $P(D\mid A)\,P(A)$;
\item \textbf{b)} oublie les poids $P(A)$ et $P(B)$ au numérateur;
\item \textbf{d)} calcule $P(B\mid D)$ et non $P(A\mid D)$;
\item \textbf{e)} utilise $P(A)=0{,}6$ dans les deux termes du dénominateur.
\end{itemize}
\end{solution}

\vspace{0.2in}

\question[3]
Un centre d'appels reçoit des appels selon un processus de Poisson de taux $4$ par heure. On s'intéresse au temps d'attente jusqu'au prochain appel. Quelle est l'expression correcte pour calculer la probabilité que ce temps soit \emph{supérieur} à $30$ minutes?\\

\begin{checkboxes}
\choice $P(T>0{,}5)$, où $T\sim \mathrm{Exp}(8)$
\correctchoice $P(T>0{,}5)$, où $T\sim \mathrm{Exp}(4)$
\choice $P(X>2)$, où $X\sim \mathrm{Poisson}(4)$
\choice $P(X>0{,}5)$, où $X\sim \mathrm{Poisson}(4)$
\choice $P(T>0{,}5)$, où $T\sim \mathrm{Gamma}(4,4)$
\end{checkboxes}

\begin{solution}
La bonne réponse est \textbf{b)}.

Dans un processus de Poisson de taux $4$ par heure, le \textbf{temps d'attente jusqu'au prochain événement} suit une loi exponentielle de paramètre $4$. On s'intéresse à un temps supérieur à $0{,}5$ heure (soit 30 minutes), d'où $P(T>0{,}5)$ avec $T\sim\mathrm{Exp}(4)$.

Problèmes avec les autres réponses :
\begin{itemize}
\item \textbf{a)} Le paramètre $8$ ne correspond pas au taux du processus de Poisson;
\item \textbf{c)} et \textbf{d)} utilisent une loi de Poisson, qui modélise un \emph{nombre} d'événements, pas un temps d'attente;
\item \textbf{e)} La loi Gamma conviendrait au temps d'attente jusqu'au $k$-ième appel, pas jusqu'au prochain.
\end{itemize}
\end{solution}

\vspace{0.2in}

\question Dites si les énoncés suivants sont vrais ou faux. Aucune justification n'est requise.

\begin{parts}

    \part[1] Si $X \sim \mathrm{Bin}(n, p)$, alors $E(X) = np$ et $\mathrm{Var}(X) = np(1-p)$.

    \begin{oneparcheckboxes}
        \correctchoice Vrai
        \choice Faux
    \end{oneparcheckboxes}

    \begin{solution}
    Vrai. Ce sont les formules classiques de l'espérance et de la variance de la loi binomiale.
    \end{solution}

    \part[1] Si $X \sim \mathrm{Exp}(\lambda)$, alors $\mathrm{Var}(X) = \lambda^2$.

    \begin{oneparcheckboxes}
        \choice Vrai
        \correctchoice Faux
    \end{oneparcheckboxes}

    \begin{solution}
    Faux. Pour une loi exponentielle de paramètre $\lambda$, on a $\mathrm{Var}(X)=\dfrac{1}{\lambda^2}$ (et $E(X)=\dfrac{1}{\lambda}$).
    \end{solution}

\end{parts}

\vspace{0.5cm}

\question[2]
Soit $X \sim \mathcal{N}(4,\,9)$ et $Y \sim \mathcal{N}(2,\,1)$, deux variables aléatoires indépendantes. Quelle est la distribution de $W = 3X - 2Y$?\\

\begin{checkboxes}
\choice $\mathcal{N}(8,\,10)$\\[4pt]
\choice $\mathcal{N}(8,\,29)$\\[4pt]
\correctchoice $\mathcal{N}(8,\,85)$\\[4pt]
\choice $\mathcal{N}(16,\,85)$\\[4pt]
\choice $\mathcal{N}(8,\,130)$
\end{checkboxes}

\begin{solution}
La bonne réponse est \textbf{c)}.
\begin{align*}
E(W)&=3E(X)-2E(Y)=3(4)-2(2)=8.\\
\mathrm{Var}(W)&=3^2\mathrm{Var}(X)+2^2\mathrm{Var}(Y)=9(9)+4(1)=81+4=85.
\end{align*}
Donc $W\sim\mathcal{N}(8,\,85)$.

Problèmes avec les autres réponses :
\begin{itemize}
\item \textbf{a)} additionne simplement les variances sans tenir compte des coefficients;
\item \textbf{b)} multiplie les variances par les coefficients (sans les mettre au carré) : $3(9)+2(1)=29$;
\item \textbf{d)} utilise $3E(X)+2E(Y)=16$ pour la moyenne au lieu de $3E(X)-2E(Y)=8$;
\item \textbf{e)} multiplie $(3^2+2^2)(9+1)=130$.
\end{itemize}
\end{solution}

\newpage

%%%%%%%%%%%%%%%% STATISTIQUES DESCRIPTIVES %%%%%%%%%%%%%%%%

\question[3]

Supposons le petit échantillon suivant :
\[
2,\;4,\;4,\;6,\;9.
\]

Ajoutez une observation à ce jeu de données de façon à diminuer la moyenne sans modifier la médiane. Justifiez brièvement votre réponse.

\begin{solution}
Une réponse possible est d'ajouter \(1\).

Le nouveau jeu de données ordonné est
\[
1,\;2,\;4,\;4,\;6,\;9.
\]

La médiane reste égale à
\[
\frac{4+4}{2}=4.
\]

La moyenne initiale est
\[
\bar{x}=\frac{2+4+4+6+9}{5}=5.
\]

La nouvelle moyenne est
\[
\bar{x}_{\text{nouvelle}}=\frac{1+2+4+4+6+9}{6}
=
\frac{26}{6}
\approx 4{,}33.
\]

La moyenne a donc diminué.
\end{solution}

\vspace{2in}

\question[6]

Vous voulez connaître l'âge des gens présents à une fête de quartier. Un ami questionne tous les individus présents et vous remet les résultats sous la forme suivante: 


\begin{table}[h!]
\begin{center}
\begin{tabular}{c|c}
Classe d'âge & Fréquence \\
\hline
$[0,5]$ & 8 \\
$[5,10]$ & 15 \\
$[10, 25]$ & 6 \\
$[25, 40]$ & ? \\
$[40, 100]$ & 10 
\end{tabular}
\end{center}
\end{table}

Le nombre d'individus âgés entre 25 et 40 ans est illisible. Votre ami avait toutefois pris le temps de faire un histogramme des données et vous informe que la hauteur du rectangle pour le groupe des individus âgés de 5 à 10 ans est de 0,06. Déduisez le nombre total de personnes présentes à la fête de quartier. 

\begin{solution}
Dans un histogramme, la hauteur $h$ d'une classe est donnée par: 
\begin{align*}
h=\frac{\mbox{Fréquence relative}}{\mbox{longueur de la classe}}
\end{align*}
donc
\begin{align*}
\mbox{Fréquence relative}=hauteur\times\mbox{longueur de la classe}=0.06\times(10-5)=0.3
\end{align*}
Aussi, la fréquence relative est donnée par: 
\begin{align*}
\mbox{Fréquence relative}=\frac{\mbox{Fréquence}}{\mbox{nombre total d'observations}}
\end{align*}
On déduit donc que: 
\begin{align*}
\mbox{nombre total d'observations} =\frac{\mbox{Fréquence}}{\mbox{Fréquence relative}}=\frac{15}{0.3}=50
\end{align*}
Le nombre total de personnes présentes à la fête de quartier est donc de $n=50$.    
\end{solution}

\newpage


%%%%%%%%%%%%%%%% LOIS ÉCHANTILLONNALES %%%%%%%%%%%%%%%%

\question[6]
Soit \(X_1,\ldots,X_{17}\) un échantillon aléatoire simple issu d'une distribution $\mathcal{N}(\mu=5,\sigma^2=4)$,
et soit \(S^2\) la variance échantillonnale.

Trouvez l'espérance de la variable
\[
Y=4S^2.
\]
\vspace{2.5in}
\begin{solution}
Comme \(S^2\) est la variance échantillonnale, on sait que
\[
E(S^2)=\sigma^2.
\]

Ici,
\[
\sigma^2=4.
\]

Donc,
\[
E(S^2)=4.
\]

On cherche
\[
E(Y)=E(4S^2).
\]

Par linéarité de l'espérance,
\[
E(Y)=4E(S^2).
\]

Donc,
\[
E(Y)=4(4)=16.
\]

Ainsi,
\[
\boxed{E(Y)=16.}
\]
\end{solution}

%%%%%%%%%%%%%%%% INTERVALLES DE CONFIANCE %%%%%%%%%%%%%%%%
\newpage

\question[10]
On a construit un intervalle de confiance pour la proportion de personnes qui préfèrent le chocolat noir au chocolat au lait. On a utilisé un échantillon de \(2400\) personnes, parmi lesquelles \(960\) ont répondu qu'elles préféraient le chocolat noir.

On a obtenu l'intervalle
\[
[37{,}83\%;\ 42{,}17\%].
\]

Quel est le niveau de confiance de cet intervalle?

\begin{solution}
On a
\[
\hat{p}=\frac{960}{2400}=0{,}40.
\]

L'intervalle obtenu est
\[
[0{,}3783;\ 0{,}4217].
\]

La marge d'erreur est donc
\[
0{,}4217-0{,}40=0{,}0217.
\]

Pour un intervalle de confiance pour une proportion, la forme de l'intervalle est
\[
\hat{p}\pm z_{\alpha/2}\sqrt{\frac{\hat{p}(1-\hat{p})}{n}}.
\]

Ici,
\[
\sqrt{\frac{\hat{p}(1-\hat{p})}{n}}
=
\sqrt{\frac{0{,}40(0{,}60)}{2400}}
=
\sqrt{0{,}0001}
=
0{,}01.
\]

On cherche donc la valeur de \(z_{\alpha/2}\) telle que
\[
z_{\alpha/2}(0{,}01)=0{,}0217.
\]

Ainsi,
\[
z_{\alpha/2}=2{,}17.
\]

Le niveau de confiance est donc
\[
P(-2{,}17\leq Z\leq 2{,}17),
\]
où \(Z\sim N(0,1)\).

À l'aide de la table de la loi normale,
\[
P(-2{,}17\leq Z\leq 2{,}17)
=
2\Phi(2{,}17)-1.
\]

Comme
\[
\Phi(2{,}17)\approx 0{,}9850,
\]
on obtient
\[
2(0{,}9850)-1=0{,}9700.
\]

Le niveau de confiance est donc environ
\[
\boxed{97\%.}
\]
\end{solution}

%%%%%%%%%%%%%%%% INTERVALLES DE CONFIANCE À DEUX POPULATIONS %%%%%%%%%%%%%%%%
\newpage
\question[10]
Deux entreprises produisent des batteries rechargeables. On s'intéresse à la durée de vie moyenne, en heures, de leurs batteries.

Pour l'entreprise A, on prélève un échantillon de taille \(n_A=36\) et on obtient \(\bar x_A=82\). On sait que l'écart-type de la population vaut \(\sigma_A=12\).

Pour l'entreprise B, on prélève un échantillon de taille \(n_B=49\) et on obtient \(\bar x_B=76\). On sait que l'écart-type de la population vaut \(\sigma_B=14\).

Construisez un intervalle de confiance à \(95\%\) pour la différence des moyennes \(\mu_A-\mu_B\).

\begin{solution}
Comme les variances des deux populations sont connues, on utilise l'intervalle de confiance
\[
(\bar x_A-\bar x_B)\pm z_{\alpha/2}\sqrt{\frac{\sigma_A^2}{n_A}+\frac{\sigma_B^2}{n_B}}.
\]

On a
\[
\bar x_A-\bar x_B=82-76=6.
\]

L'erreur-type vaut
\[
\sqrt{\frac{12^2}{36}+\frac{14^2}{49}}
=\sqrt{4+4}
=\sqrt{8}
\approx 2{,}828.
\]

Pour un niveau de confiance de \(95\%\), on utilise \(z_{0{,}025}=1{,}96\).

La marge d'erreur est donc
\[
1{,}96\times 2{,}828\approx 5{,}54.
\]

Ainsi, l'intervalle de confiance est
\[
6\pm 5{,}54,
\]
soit
\[
[0{,}46;\ 11{,}54].
\]

Donc, un intervalle de confiance à \(95\%\) pour \(\mu_A-\mu_B\) est
\[
\boxed{[0{,}46;\ 11{,}54]}.
\]
\end{solution}

\vspace{2in}



%%%%%%%%%%%%%%%% TESTS D'HYPOTHÈSES %%%%%%%%%%%%%%%%

\newpage 

\question[8]
Je veux tester les hypothèses suivantes en pigeant un échantillon de taille \(n=25\) dans une population normale de variance inconnue :
\[
H_0:\mu=21
\qquad \text{contre} \qquad
H_1:\mu<21.
\]

J'obtiens un seuil observé de \(0{,}8413\).

Quel serait le seuil observé pour le test bilatéral avec la même hypothèse nulle?

\begin{solution}
Le seuil observé donné correspond au test unilatéral à gauche :
\[
H_1:\mu<21.
\]

Ici, le seuil observé unilatéral est
\[
0{,}8413.
\]

Cela signifie que la statistique observée n'est pas dans la queue gauche de la distribution. Elle est plutôt du côté opposé à l'hypothèse alternative.

Pour le test bilatéral correspondant,
\[
H_0:\mu=21
\qquad \text{contre} \qquad
H_1:\mu\neq 21,
\]
le seuil observé bilatéral est
\[
2\min(0{,}8413,\;1-0{,}8413).
\]

Donc,
\[
2\min(0{,}8413,\;0{,}1587)
=
2(0{,}1587)
=
0{,}3174.
\]

Ainsi, le seuil observé pour le test bilatéral correspondant est
\[
\boxed{0{,}3174}.
\]
\end{solution}

\newpage

\question[10]
Deux chaînes d'assemblage produisent des composantes électroniques. On veut vérifier si la chaîne A produit en moyenne des composantes plus lourdes que la chaîne B.

On prélève deux échantillons indépendants et on obtient les résultats suivants :
\begin{itemize}
\item Chaîne A : \(n_A=64\), \(\bar x_A=503\) mg, \(\sigma_A=8\) mg;
\item Chaîne B : \(n_B=100\), \(\bar x_B=500\) mg, \(\sigma_B=10\) mg.
\end{itemize}

On suppose que les deux populations sont normales et que les variances sont connues.

Au seuil \(\alpha=5\%\), testez
\[
H_0:\mu_A-\mu_B=0
\qquad \text{contre} \qquad
H_1:\mu_A-\mu_B>0.
\]

\begin{solution}
Comme les variances des deux populations sont connues, on utilise un test \(Z\) à deux échantillons.

La statistique de test est
\[
Z_0=\frac{(\bar x_A-\bar x_B)-0}{\sqrt{\sigma_A^2/n_A+\sigma_B^2/n_B}}.
\]

On a
\[
\bar x_A-\bar x_B=503-500=3
\]
et
\[
\sqrt{\frac{8^2}{64}+\frac{10^2}{100}}
=\sqrt{1+1}
=\sqrt{2}
\approx 1{,}414.
\]

Donc,
\[
Z_0=\frac{3}{1{,}414}\approx 2{,}12.
\]

Au seuil \(\alpha=5\%\), pour un test unilatéral à droite, on rejette \(H_0\) si
\[
Z_0>z_{0{,}05}=1{,}645.
\]

Comme
\[
2{,}12>1{,}645,
\]
on rejette \(H_0\) au seuil de \(5\%\).

On conclut qu'il y a une évidence statistique que la chaîne A produit en moyenne des composantes plus lourdes que la chaîne B.
\end{solution}
\vspace{2.5in}

\newpage
\question[]

Caroline a trouvé une pièce de monnaie. Son intuition lui dit que la pièce tombe plus souvent sur le côté PILE que sur le côté FACE. Elle souhaite vérifier cette hypothèse et lance donc la pièce à \(64\) reprises dans les airs en notant le résultat obtenu.

\begin{parts}

\part[9] Au seuil de \(5\%\), quelle est la puissance de son test statistique si, en vérité, la pièce de monnaie a une probabilité de \(70\%\) de tomber sur le côté PILE?

\begin{solution}
Soit
\[
X=\text{le nombre de fois où la pièce tombe sur PILE parmi les 64 lancers}.
\]

On veut tester
\[
H_0:p=0{,}5
\qquad \text{contre} \qquad
H_1:p>0{,}5.
\]

Sous \(H_0\), on a
\[
X\sim \mathrm{Bin}(64,0{,}5).
\]

Comme l'hypothèse alternative est unilatérale à droite, la région de rejet est de la forme
\[
X\geq c.
\]

On cherche donc une valeur de \(c\) telle que
\[
P_{0{,}5}(X\geq c)\leq 0{,}05.
\]

On peut vérifier que
\[
P_{0{,}5}(X\geq 39)\approx 0{,}0517
\]
et que
\[
P_{0{,}5}(X\geq 40)\approx 0{,}0300.
\]

Ainsi, au seuil de \(5\%\), on rejette \(H_0\) si
\[
X\geq 40.
\]

On cherche maintenant la puissance du test lorsque la vraie probabilité d'obtenir PILE est \(p=0{,}7\).

Sous cette valeur de \(p\), on a
\[
X\sim \mathrm{Bin}(64,0{,}7).
\]

La puissance est donc
\[
P_{0{,}7}(X\geq 40).
\]

Ainsi,
\[
P_{0{,}7}(X\geq 40)
=
\sum_{x=40}^{64}
\binom{64}{x}
(0{,}7)^x(0{,}3)^{64-x}.
\]

Numériquement,
\[
P_{0{,}7}(X\geq 40)\approx 0{,}9236.
\]

La puissance du test est donc environ
\[
\boxed{0{,}924}.
\]

Le test a donc environ \(92{,}4\%\) de chances de rejeter \(H_0\) si la vraie probabilité d'obtenir PILE est \(0{,}7\).
\end{solution}
\vfill
\part[3] Caroline avait aussi pensé utiliser un seuil de \(10\%\). Cochez la bonne réponse ci-dessous.

\begin{checkboxes}
\choice La puissance aurait été plus petite avec un seuil de \(10\%\).
\correctchoice La puissance aurait été plus grande avec un seuil de \(10\%\).
\choice La puissance aurait été identique avec un seuil de \(10\%\).
\choice On ne peut pas savoir avec l'information donnée si la puissance aurait été plus grande, plus petite ou identique avec un seuil de \(10\%\).
\end{checkboxes}

\begin{solution}
En augmentant le seuil du test de \(5\%\) à \(10\%\), on augmente la probabilité de rejeter \(H_0\).

La région de rejet devient donc plus grande ou reste la même. La puissance, qui est la probabilité de rejeter \(H_0\) lorsque l'hypothèse alternative est vraie, ne peut donc pas diminuer.

Dans ce contexte, la puissance serait plus grande avec un seuil de \(10\%\).
\end{solution}

\end{parts}

\newpage

%%%%%%%%%%%%%%%% TESTS APPARIÉS %%%%%%%%%%%%%%%%

\question[]
On veut vérifier si un programme d'entraînement physique de 6 semaines réduit la pression artérielle systolique (en mmHg) d'adultes sédentaires. Six participants sont sélectionnés au hasard et leur pression artérielle est mesurée avant et après le programme.

\begin{center}
\renewcommand{\arraystretch}{1.5}
\begin{tabular}{|l|cccccc|c|}
\hline
\multirow{2}{*}{Mesure} & \multicolumn{6}{c|}{Participant} & \multirow{2}{*}{Moyenne} \\
\cline{2-7}
 & 1 & 2 & 3 & 4 & 5 & 6 & \\
\hline
Avant & 152 & 145 & 163 & 148 & 155 & 149 & 152 \\
Apr\`es & 146 & 140 & 155 & 143 & 151 & 147 & 147 \\
\hline
\end{tabular}
\end{center}

On pose $d_i = \text{Avant}_i - \text{Après}_i$ et on vous informe que:
\[
s_D^2 = 4 \qquad \text{et} \qquad \bar{d} = 5.
\]

\begin{parts}

\part[5] Au seuil $\alpha=5\%$, testez si le programme réduit la pression artérielle en moyenne. 

\begin{solution}
On pose $\mu_D = E(\text{Avant}_i - \text{Après}_i)$ et on teste
\[
H_0:\mu_D=0 \qquad\text{contre}\qquad H_1:\mu_D>0.
\]

La statistique de test est
\[
T_0=\frac{\bar{d}}{s_D/\sqrt{n}}=\frac{5}{2/\sqrt{6}}=\frac{5\sqrt{6}}{2}\approx 6{,}12.
\]

Sous $H_0$, $T_0\sim t_5$. La règle de décision est de rejeter $H_0$ si $T_0>t_{0{,}05;\,5}=2{,}015$.

Comme $6{,}12>2{,}015$, on rejette $H_0$ au seuil $5\%$.

On conclut qu'il y a une évidence statistique que le programme réduit la pression artérielle systolique en moyenne.
\end{solution}

\vfill

\part[2] Pourquoi est-il probablement préférable d'utiliser des échantillons appariés dans ce contexte~? Une seule réponse est bonne.

\begin{checkboxes}
\choice Car tous les participants ont naturellement la même pression artérielle de base.
\choice Car la mesure après le programme est plus précise que la mesure avant.
\correctchoice Car nous voulons nous protéger contre la possibilité que les différences naturelles de pression artérielle entre les participants influencent la décision du test.
\choice Car nous rejetons $H_0$ au seuil $\alpha=5\%$.
\choice Car il s'agit d'observations prises dans le temps.
\end{checkboxes}

\begin{solution}
En appariant les mesures avant et après pour chaque participant, on élimine la variabilité due aux différences entre individus (certains ont naturellement une pression plus élevée que d'autres). Le test porte alors uniquement sur la différence intra-individu attribuable au programme.
\end{solution}

\end{parts}

\newpage

%%%%%%%%%%%%%%%% RLS %%%%%%%%%%%%%%%%


\question[]

On a ajusté un modèle de régression linéaire pour le prix de vente d'une maison, en milliers de dollars, en fonction de l'évaluation de cette maison, aussi en milliers de dollars, à l'aide d'un jeu de données de \(92\) observations. On a obtenu les valeurs suivantes :
\[
\hat{\beta}_0 = 20{,}942,
\qquad
\hat{\beta}_1 = 1{,}06873,
\qquad
\hat{\sigma}^2 = 1075.
\]

De plus, on vous indique que
\[
\bar{x} = 201{,}7504
\qquad \text{et} \qquad
S_{xx} = 5209571.
\]

\begin{parts}

\part[3] Interprétez la pente estimée, en contexte.

\begin{solution}
La pente estimée est
\[
\hat{\beta}_1=1{,}06873.
\]

Cela signifie que, selon le modèle ajusté, lorsqu'une maison a une évaluation plus élevée de \(1\) millier de dollars, son prix de vente moyen prévu augmente d'environ \(1{,}06873\) millier de dollars.

Autrement dit, une augmentation de \(1\,000\) \$ de l'évaluation est associée à une augmentation moyenne d'environ
\[
1\,069\ \text{\$}
\]
du prix de vente.
\end{solution}

\vspace{1.5in}

\part[6] Un ami vient de mettre en vente sa maison, qui a été évaluée à \(280\) milliers de dollars.

Construisez un intervalle de confiance à \(95\%\) pour le prix de vente de sa maison.

\begin{solution}
Comme on s'intéresse au prix de vente d'une maison individuelle, on construit un intervalle de prédiction.

Pour \(x_0=280\), la valeur prédite est
\[
\hat{y}_0
=
\hat{\beta}_0+\hat{\beta}_1x_0.
\]

Donc,
\[
\hat{y}_0
=
20{,}942+1{,}06873(280)
=
320{,}1864.
\]

L'intervalle de prévision à \(95\%\) est donné par
\[
\hat{y}_0
\pm
t_{0{,}975;\,n-2}
\hat{\sigma}
\sqrt{
1+\frac{1}{n}+\frac{(x_0-\bar{x})^2}{S_{xx}}
}.
\]

Ici, \(n=92\), donc les degrés de liberté sont
\[
n-2=90.
\]

On utilise
\[
t_{0{,}975;\,90}\approx 1{,}987.
\]

De plus,
\[
\hat{\sigma}
=
\sqrt{1075}.
\]

Donc,
\[
\hat{\sigma}
\sqrt{
1+\frac{1}{n}+\frac{(x_0-\bar{x})^2}{S_{xx}}
}
=
\sqrt{
1075
\left(
1+\frac{1}{92}
+
\frac{(280-201{,}7504)^2}{5209571}
\right)
}.
\]

Ainsi,
\[
\hat{\sigma}
\sqrt{
1+\frac{1}{n}+\frac{(x_0-\bar{x})^2}{S_{xx}}
}
\approx
32{,}984.
\]

L'intervalle est donc
\[
320{,}1864
\pm
1{,}987(32{,}984).
\]

On obtient
\[
320{,}1864
\pm
65{,}539.
\]

L'intervalle de prédiction à \(95\%\) est donc environ
\[
[254{,}647,\;385{,}726].
\]

En contexte, on prévoit que le prix de vente de cette maison se situera entre environ
\[
254\,647\ \text{\$}
\quad \text{et} \quad
385\,726\ \text{\$}.
\]
\end{solution}

\vspace{3 in}
\newpage

\part[4] Au seuil de \(5\%\), peut-on conclure qu'il existe une relation linéaire significative entre l'évaluation d'une maison et son prix de vente?

\begin{solution}
On veut tester
\[
H_0:\beta_1=0
\qquad \text{contre} \qquad
H_1:\beta_1\neq 0.
\]

La statistique de test est
\[
T=
\frac{\hat{\beta}_1-0}{\mathrm{SE}(\hat{\beta}_1)}.
\]

On a
\[
\mathrm{SE}(\hat{\beta}_1)
=
\sqrt{\frac{\hat{\sigma}^2}{S_{xx}}}.
\]

Donc,
\[
\mathrm{SE}(\hat{\beta}_1)
=
\sqrt{\frac{1075}{5209571}}
\approx
0{,}01437.
\]

Ainsi,
\[
T
=
\frac{1{,}06873}{0{,}01437}
\approx
74{,}38.
\]

Sous \(H_0\), la statistique de test suit une loi de Student à
\[
n-2=90
\]
degrés de liberté.

Au seuil de \(5\%\), pour un test bilatéral, la valeur critique est environ
\[
t_{0{,}975;90}\approx 1{,}987.
\]

Comme
\[
|74{,}38|>1{,}987,
\]
on rejette \(H_0\).

On peut donc conclure, au seuil de \(5\%\), qu'il existe une relation linéaire significative entre l'évaluation d'une maison et son prix de vente.
\end{solution}

\vspace{3.5in}

\end{parts}

\end{questions}
\end{document} 
